\chapter{Criterio de distinción operativa}

\section{Distinción disciplinada}
\begin{definition}
Dos piezas son distinguibles en sentido operativo cuando puede declararse entre ellas una diferencia estructural, jerárquica o funcional que afecte de manera real su tratamiento dentro del canon.
\end{definition}

\begin{observation}
La distinción operativa no depende de variación verbal superficial. Exige una diferencia que modifique el modo de lectura, comparación, transición o uso disciplinado de las piezas.
\end{observation}

\section{Condición mínima de distinción}
\begin{definition}
Existe distinción operativa mínima cuando al menos una de las siguientes diferencias puede ser declarada sin ambigüedad:
\begin{enumerate}
  \item diferencia de dependencia canónica;
  \item diferencia de nivel jerárquico;
  \item diferencia de umbral operativo;
  \item diferencia de función dentro del corpus;
  \item diferencia de información conservada bajo lectura primitiva o proyectada.
\end{enumerate}
\end{definition}

\begin{observation}
No toda diferencia observable alcanza rango operativo. Debe ser relevante para el tratamiento disciplinado y no solo para descripción incidental.
\end{observation}

\section{Distinción y comparabilidad}
\begin{definition}
La comparabilidad no elimina la distinción. Dos piezas comparables pueden seguir siendo operativamente distintas si la relación entre ellas conserva diferencia real de función, nivel o estructura.
\end{definition}

\begin{observation}
Esta cláusula evita confundir comparabilidad con homogeneidad.
\end{observation}

\section{Distinción y transición}
\begin{definition}
Cuando una formulación fortalece a otra mediante transición disciplinada, la distinción operativa entre ambas debe permanecer legible y no quedar absorbida por mera reformulación nominal.
\end{definition}

\begin{observation}
Si una transición vuelve indistinguibles sus extremos, la ganancia operativa queda debilitada o resulta dudosa.
\end{observation}

\section{Condición mínima de fallo}
\begin{definition}
Se considera fallida una distinción operativa cuando la diferencia alegada entre dos piezas no produce variación discernible en su dependencia, su nivel, su umbral o su función.
\end{definition}

\begin{observation}
Esta cláusula protege el canon intermedio frente a proliferación de categorías o bloques que solo cambian de nombre.
\end{observation}

\section{Cláusula de alcance}
\begin{observation}
El criterio de distinción operativa fija una regla mínima para separar piezas del corpus sin colapsarlas en continuidad indiferenciada. No constituye todavía una teoría completa de clasificación.
\end{observation}

